\section{Methods description form (Track 1 methods)}
\textit{Answers to the organisers' official methods description template (version with the required LLM question), reproduced here because the submission form has no separate upload field for it.}

\begin{description}[leftmargin=0pt,style=nextline,itemsep=4pt,font=\sffamily\bfseries\small\color{branddark}]
\item[Team name] aperezcastro (Agustín Pérez Castro, individual)
\item[Model number (fill out this form once per model/approach submitted; up to 6)] 1 (BUB1B compound heterozygous; aperezcastro\_bub1b\_comphet.csv)
\item[Please describe your model/approach in detail.] Three-arm pipeline on the provided VCF, run as time-capped jobs on Modal (scripts in the GitHub repo). (1) bcftools norm (split multiallelics, left-align vs Ensembl GRCh38) $\rightarrow$ Ensembl VEP 116 offline (MANE, HGVS, gnomAD v4.1 exomes+genomes, ClinVar with review status, SIFT, PolyPhen-2, AlphaMissense). (2) Unbiased genome-wide funnel: PASS $\rightarrow$ max gnomAD AF < 1\% $\rightarrow$ HIGH/MODERATE/splice $\rightarrow$ gene-level recessive model ($\geq$2 het or 1 hom) = 339 variants in 99 genes, each reviewed. (3) Hypothesis-driven arm: MVA/spindle-assembly-checkpoint panel (BUB1B, CEP57, TRIP13, CENATAC, BUB1, BUB3, MAD1L1, MAD2L1BP, SMC5/6 complex...) at any consequence and AF < 5\%; of the 42 rare PASS variants, the 32 inside annotated gene bodies were scored with SpliceAI 1.3.1 (50 and 2000 nt windows; all max delta $\leq$ 0.03). (4) Exomiser 15.1 (data 2512) with the 8 HPO terms, with and without a phenotype pre-filter: BUB1B ranks \#1 genome-wide in both. (5) ACMG/AMP classification with ClinGen SVI refinements and Bayesian points (Tavtigian 2020) for each submitted variant; phase feasibility tested directly (only 3 heterozygous sites in the 10.9 kb interval, gaps 6.8 and 4.1 kb: not phasable with short reads) and argued from biology. (6) Genome-level QC from aggregate statistics: allelic balance per autosome and in 5-Mb bins (excess overdispersion confined to pericentromeric/segmental-duplication regions, i.e. technical; no detectable whole-chromosome imbalance, as expected for variegated low-level mosaicism), balanced allele fractions around BUB1B, and no run of homozygosity > 2 Mb (consistent with compound heterozygosity rather than autozygosity). Answer: BUB1B c.2210T>G p.(Leu737*) [Likely pathogenic, 9 points; Pathogenic with ClinVar corroboration] + c.3006T>G p.(Asn1002Lys) [VUS, 4-5 points], MVA type 1.
\item[(required) If commercially-available Generative AI was used, record the provider, the plan or tier, and the relevant setting in your methods description. A line is enough, for example, "Anthropic API, Claude Sonnet, commercial terms, no training on customer content."] Anthropic Claude via the Claude Code CLI (Claude Opus models), Claude Max (consumer tier), model-training opt-out enabled ('Help improve Claude' off), no feedback provided on any output, no content flagged for safety review - per the organisers' confirmed guidance of 28 Sep 2026 (discussion \#2). Applicable terms: opted-out conversations are not used for training and are retained 30 days. Additional data minimisation: the model only saw code, aggregate counts, HPO terms and short candidate lists ($\leq$ 50 annotated variants); raw data and genome-scale derivatives were blocked from the agent by configuration. No patient variant was sent to any third-party API.
\item[Please state whether the submission file is the automated output of your computational approach (preferred), or if it has undergone downstream manual review and curation.] Candidate generation is automated (funnel, panel, SpliceAI, Exomiser); the final 3-row file and EPCR values were curated manually from the automated evidence.
\item[Please describe any downstream manual review in detail.] Manual review of the 99 recessive candidate genes (artefact-prone regions, ClinVar-benign, low depth, phenotype fit), of Exomiser's top hits, and of dominant rhabdomyosarcoma predisposition genes. Row 1 = the pair (EPCR 0.95). Rows 2-3 are low-probability alternative pairings with a third rare deep-intronic BUB1B variant (SpliceAI 0.00; EPCR 0.03 and 0.01). EPCR reflects residual uncertainty about phase and unobserved SV/regulatory alleles. Validated with the organisers' evaluation.py.
\item[Please state whether your approach only used publicly available data (preferred) or if proprietary data was also used. If proprietary data were used, where possible, please also include a submission based only on publicly available data.] Only publicly available data.
\item[Please describe any public data used in detail (e.g., reference panels, population databases, literature).] GRCh38 (Ensembl 116 primary assembly), Ensembl VEP 116 cache, GENCODE (SpliceAI grch38), gnomAD v4.1, ClinVar (NCBI VCF, Oct 2026), AlphaMissense hg38, Exomiser data release 2512 (REVEL, MVP, AlphaMissense, SpliceAI, HPO/MGI/ZFIN phenotypes), Human Phenotype Ontology (hp.obo), UniProt O60566 and the AlphaFold model, PubMed literature (cited in the report).
\item[Please describe any proprietary data used in detail.] None.
\item[Is your approach able to output proposed pairs of compound heterozygous candidate variants, or only single candidate variants?] Yes: pairs are produced by the gene-level recessive model (all qualifying pairs within a gene) and ranked; single variants are produced for dominant models.
\item[How did you handle secondary or incidental findings, if any? (See the finding\_type column in your submission.)] No secondary findings were submitted. Incidental findings were handled following the organisers' IRB guidance (discussion \#20) and are not reported.
\item[Please provide an estimate of run time and cost for your approach (if possible).] About 1 hour wall-clock of cloud compute (annotation parallelised over 25 containers \textasciitilde{}20 min, Exomiser \textasciitilde{}7 min, SpliceAI \textasciitilde{}4 min) plus reference downloads; total Modal cost \textasciitilde{}1.4 USD (scripts/modal/gastos.csv). Reproducible from the repository.
\item[Provide a method abstract (up to 500 words). Please include strengths and limitations of your method as well as methodological detail.] We combined an unbiased genome-wide search with a hypothesis-driven analysis of genes linked to mosaic variegated aneuploidy (MVA). The single-sample WGS VCF was normalised and annotated offline with Ensembl VEP 116 (gnomAD v4.1, ClinVar, AlphaMissense) on a private cloud environment. A recessive gene-level funnel reduced 4.96 M variants to 339 variants in 99 genes, of which only BUB1B belongs to the MVA/spindle-checkpoint panel. Because the second, hypomorphic allele in MVA1 is often subtle, the panel was searched at every consequence class and scored with SpliceAI; independently, Exomiser 15.1 with the patient's HPO terms ranked BUB1B first genome-wide with and without a phenotype filter (score 0.59 dominant / 0.56 recessive vs $\leq$ 0.22 for the next gene). The proband carries BUB1B c.2210T>G p.(Leu737*) (nonsense, NMD predicted; ClinVar P/LP; ACMG 9 points, Likely pathogenic, Pathogenic with ClinVar corroboration) and c.3006T>G p.(Asn1002Lys) (ultra-rare kinase-domain missense, AlphaMissense 0.92, buried high-confidence residue in the AlphaFold model 12 Å from the known MVA1 allele L1012P; 4-5 points, VUS; its class, low abundance or impaired function, is unmeasured). This null + kinase-domain hypomorph architecture is the classical MVA1 pattern and mirrors a published mouse model (BubR1 X753 + L1002P, mouse L1002 = human L1012). A third rare deep-intronic BUB1B variant was assessed and deprioritised (SpliceAI 0.00). Phase could not be observed: the 10.9 kb interval contains only three heterozygous sites (gaps 6.8 and 4.1 kb), so we tested and documented that short-read phasing is infeasible and argued trans configuration from biology (monoallelic loss does not cause MVA), absence of any other qualifying variant and balanced allele fractions, and the genome shows no run of homozygosity > 2 Mb. Strengths: three independent arms converging on the same pair, explicit ACMG evidence with Bayesian points, sensitivity to subtle alleles, full reproducibility with cost-capped scripts, and strict data minimisation (no third-party variant APIs). Limitations: no BAM/CRAM, so structural variants were not called genome-wide; phase is inferred rather than observed; in silico predictors disagree in magnitude for the missense (REVEL 0.47); single proband without segregation data.
\end{description}
